% =============================================================
% intro_generale.tex — Introduction générale (~2 pages visées)
% Règle du guide : poser la problématique SANS révéler de résultat.
% À finaliser en dernier, une fois les chapitres 1 à 5 stabilisés.
% =============================================================
\chapter*{Introduction générale}
\addcontentsline{toc}{chapter}{Introduction générale}
\thispagestyle{plain}
\label{ch:intro}

\begin{attentedoc}
Squelette d'introduction. Chaque paragraphe ci-dessous doit être réécrit une fois le
chapitre~1 (problématique, \S\ref{sec:problematique}) et les chapitres~3 à 5 stabilisés,
pour que l'introduction annonce fidèlement un plan définitif.
\end{attentedoc}

\todo{Paragraphe de contexte général : dépendance croissante de MENAL aux applications
hébergées, exigences de sécurité et de traçabilité qui en découlent.}

\todo{Paragraphe de problématique : reprendre \emph{textuellement} la formulation validée
au chapitre~1, \S\ref{sec:problematique} --- ne pas la reformuler ici pour éviter toute
divergence entre les deux versions.}

\todo{Paragraphe d'objectifs : reprendre la liste d'objectifs du chapitre~1 sous forme de
phrase continue (pas de liste à puces dans l'introduction générale).}

\todo{Paragraphe d'annonce du plan : un paragraphe par chapitre (2 à 5), sans anticiper les
résultats du chapitre~5.}
